\documentclass[11pt, a4paper]{article}
% ── Shared preamble for all RTOS lab documents ────────────────────────────────
% Usage: % ── Shared preamble for all RTOS lab documents ────────────────────────────────
% Usage: % ── Shared preamble for all RTOS lab documents ────────────────────────────────
% Usage: \input{../preamble}  (from each labNN/ directory)
% ──────────────────────────────────────────────────────────────────────────────

% ── Packages ──────────────────────────────────────────────────────────────────
\usepackage[a4paper, top=2.5cm, bottom=2.5cm, left=2.5cm, right=2.5cm, headheight=28pt]{geometry}
\usepackage[utf8]{inputenc}
\usepackage[T1]{fontenc}
\usepackage{lmodern}
\usepackage{booktabs}
\usepackage{array}
\usepackage{tabularx}
\usepackage{longtable}
\usepackage{multirow}
\usepackage{colortbl}
\usepackage{xcolor}
\usepackage{titlesec}
\usepackage{enumitem}
\usepackage{hyperref}
\usepackage{fancyhdr}
\usepackage{parskip}
\usepackage{listings}
\usepackage[most]{tcolorbox}
\usepackage{fontawesome5}
\usepackage{microtype}
\usepackage{graphicx}
\usepackage{amsmath}

% ── Colors (KMUTNB brand) ──────────────────────────────────────────────────────
\definecolor{kmutnbBlue}{RGB}{0,56,116}
\definecolor{kmutnbGold}{RGB}{186,146,36}
\definecolor{lightGray}{RGB}{245,245,245}
\definecolor{midGray}{RGB}{200,200,200}
\definecolor{codeGray}{RGB}{248,248,248}
\definecolor{codeFrame}{RGB}{220,220,220}
\definecolor{tipteal}{RGB}{0,105,92}
\definecolor{warnAmber}{RGB}{121,85,0}
\definecolor{checkGreen}{RGB}{27,94,32}

% ── Section formatting ─────────────────────────────────────────────────────────
\titleformat{\section}
  {\large\bfseries\color{kmutnbBlue}}
  {\thesection.}{0.5em}{}
  [\color{kmutnbGold}\titlerule]

\titleformat{\subsection}
  {\normalsize\bfseries\color{kmutnbBlue}}
  {\thesubsection.}{0.5em}{}

\titleformat{\subsubsection}
  {\small\bfseries\color{kmutnbBlue}}
  {\thesubsubsection.}{0.5em}{}

\titlespacing*{\section}{0pt}{1.4ex plus .2ex}{0.6ex plus .1ex}
\titlespacing*{\subsection}{0pt}{1.0ex plus .2ex}{0.4ex plus .1ex}

% ── Listings (C and bash) ──────────────────────────────────────────────────────
\lstdefinestyle{cstyle}{
  language        = C,
  basicstyle      = \ttfamily\small,
  keywordstyle    = \bfseries\color{kmutnbBlue},
  commentstyle    = \itshape\color{gray},
  stringstyle     = \color{tipteal},
  numberstyle     = \tiny\color{midGray},
  numbers         = left,
  numbersep       = 6pt,
  stepnumber      = 1,
  backgroundcolor = \color{codeGray},
  frame           = single,
  framerule       = 0.4pt,
  rulecolor       = \color{codeFrame},
  breaklines      = true,
  showstringspaces= false,
  tabsize         = 4,
  xleftmargin     = 1.5em,
  xrightmargin    = 0.5em,
  aboveskip       = 0.8em,
  belowskip       = 0.8em,
}

\lstdefinestyle{bashstyle}{
  language        = bash,
  basicstyle      = \ttfamily\small,
  keywordstyle    = \color{kmutnbBlue},
  commentstyle    = \itshape\color{gray},
  backgroundcolor = \color{black!5},
  frame           = single,
  framerule       = 0.4pt,
  rulecolor       = \color{codeFrame},
  breaklines      = true,
  showstringspaces= false,
  xleftmargin     = 1.5em,
  xrightmargin    = 0.5em,
  aboveskip       = 0.6em,
  belowskip       = 0.6em,
}

\lstdefinestyle{textstyle}{
  basicstyle      = \ttfamily\footnotesize,
  backgroundcolor = \color{black!4},
  frame           = single,
  framerule       = 0.4pt,
  rulecolor       = \color{codeFrame},
  breaklines      = true,
  xleftmargin     = 1.5em,
  xrightmargin    = 0.5em,
  aboveskip       = 0.6em,
  belowskip       = 0.6em,
}

% ── Callout boxes ──────────────────────────────────────────────────────────────
\tcbuselibrary{skins, breakable}

\newtcolorbox{tipbox}[1][]{
  enhanced, breakable,
  colback=tipteal!8, colframe=tipteal, coltitle=white,
  fonttitle=\bfseries\small, title={\faLightbulb\quad Tip#1},
  left=6pt, right=6pt, top=4pt, bottom=4pt,
  boxrule=0.8pt, arc=3pt,
}

\newtcolorbox{warnbox}[1][]{
  enhanced, breakable,
  colback=warnAmber!8, colframe=kmutnbGold, coltitle=white,
  fonttitle=\bfseries\small, title={\faExclamationTriangle\quad Note#1},
  left=6pt, right=6pt, top=4pt, bottom=4pt,
  boxrule=0.8pt, arc=3pt,
}

\newtcolorbox{checkpointbox}[1][]{
  enhanced,
  colback=kmutnbBlue!6, colframe=kmutnbBlue, coltitle=white,
  fonttitle=\bfseries\small, title={\faCheckCircle\quad Checkpoint#1},
  left=6pt, right=6pt, top=4pt, bottom=4pt,
  boxrule=0.8pt, arc=3pt,
}

\newtcolorbox{questionbox}[1][]{
  enhanced, breakable,
  colback=kmutnbGold!10, colframe=kmutnbGold, coltitle=white,
  fonttitle=\bfseries\small, title={\faQuestion\quad Question#1},
  left=6pt, right=6pt, top=4pt, bottom=4pt,
  boxrule=0.8pt, arc=3pt,
}

% ── Checkbox list ──────────────────────────────────────────────────────────────
\newlist{checklist}{itemize}{1}
\setlist[checklist]{label=\faSquare[regular], leftmargin=1.5em, noitemsep}

% ── Misc helpers ───────────────────────────────────────────────────────────────
\newcommand{\file}[1]{\texttt{\small #1}}
\newcommand{\api}[1]{\texttt{\small #1}}

% ── Title block macro ─────────────────────────────────────────────────────────
% Usage: \labtitle{03}{Aperiodic Servers \& Producer--Consumer}{2}{CLO 1 \& CLO 2}
\newcommand{\labtitle}[4]{%
  \begin{center}
    {\color{kmutnbBlue}\rule{\linewidth}{2pt}}\\[6pt]
    {\LARGE\bfseries\color{kmutnbBlue} Laboratory #1}\\[4pt]
    {\large\bfseries #2}\\[2pt]
    {\normalsize Real-Time Operating Systems \enspace\textbullet\enspace
     M.Eng.\ ECE \enspace\textbullet\enspace KMUTNB}\\[8pt]
    {\color{kmutnbGold}\rule{\linewidth}{1pt}}
  \end{center}
  \vspace{0.3cm}
  \begin{tabular}{@{}p{3.8cm} p{10cm}@{}}
    \textbf{Platform}       & NXP FRDM-MCXN236 (ARM Cortex-M33 @ 150\,MHz) \\[2pt]
    \textbf{RTOS}           & FreeRTOS v10.6 \\[2pt]
    \textbf{Toolchain}      & ARM GNU Toolchain (\texttt{arm-none-eabi-gcc}) + Make \\[2pt]
    \textbf{Estimated time} & #3 hours \\[2pt]
    \textbf{CLO mapping}    & #4 \\
  \end{tabular}
  \vspace{0.4cm}
}
  (from each labNN/ directory)
% ──────────────────────────────────────────────────────────────────────────────

% ── Packages ──────────────────────────────────────────────────────────────────
\usepackage[a4paper, top=2.5cm, bottom=2.5cm, left=2.5cm, right=2.5cm, headheight=28pt]{geometry}
\usepackage[utf8]{inputenc}
\usepackage[T1]{fontenc}
\usepackage{lmodern}
\usepackage{booktabs}
\usepackage{array}
\usepackage{tabularx}
\usepackage{longtable}
\usepackage{multirow}
\usepackage{colortbl}
\usepackage{xcolor}
\usepackage{titlesec}
\usepackage{enumitem}
\usepackage{hyperref}
\usepackage{fancyhdr}
\usepackage{parskip}
\usepackage{listings}
\usepackage[most]{tcolorbox}
\usepackage{fontawesome5}
\usepackage{microtype}
\usepackage{graphicx}
\usepackage{amsmath}

% ── Colors (KMUTNB brand) ──────────────────────────────────────────────────────
\definecolor{kmutnbBlue}{RGB}{0,56,116}
\definecolor{kmutnbGold}{RGB}{186,146,36}
\definecolor{lightGray}{RGB}{245,245,245}
\definecolor{midGray}{RGB}{200,200,200}
\definecolor{codeGray}{RGB}{248,248,248}
\definecolor{codeFrame}{RGB}{220,220,220}
\definecolor{tipteal}{RGB}{0,105,92}
\definecolor{warnAmber}{RGB}{121,85,0}
\definecolor{checkGreen}{RGB}{27,94,32}

% ── Section formatting ─────────────────────────────────────────────────────────
\titleformat{\section}
  {\large\bfseries\color{kmutnbBlue}}
  {\thesection.}{0.5em}{}
  [\color{kmutnbGold}\titlerule]

\titleformat{\subsection}
  {\normalsize\bfseries\color{kmutnbBlue}}
  {\thesubsection.}{0.5em}{}

\titleformat{\subsubsection}
  {\small\bfseries\color{kmutnbBlue}}
  {\thesubsubsection.}{0.5em}{}

\titlespacing*{\section}{0pt}{1.4ex plus .2ex}{0.6ex plus .1ex}
\titlespacing*{\subsection}{0pt}{1.0ex plus .2ex}{0.4ex plus .1ex}

% ── Listings (C and bash) ──────────────────────────────────────────────────────
\lstdefinestyle{cstyle}{
  language        = C,
  basicstyle      = \ttfamily\small,
  keywordstyle    = \bfseries\color{kmutnbBlue},
  commentstyle    = \itshape\color{gray},
  stringstyle     = \color{tipteal},
  numberstyle     = \tiny\color{midGray},
  numbers         = left,
  numbersep       = 6pt,
  stepnumber      = 1,
  backgroundcolor = \color{codeGray},
  frame           = single,
  framerule       = 0.4pt,
  rulecolor       = \color{codeFrame},
  breaklines      = true,
  showstringspaces= false,
  tabsize         = 4,
  xleftmargin     = 1.5em,
  xrightmargin    = 0.5em,
  aboveskip       = 0.8em,
  belowskip       = 0.8em,
}

\lstdefinestyle{bashstyle}{
  language        = bash,
  basicstyle      = \ttfamily\small,
  keywordstyle    = \color{kmutnbBlue},
  commentstyle    = \itshape\color{gray},
  backgroundcolor = \color{black!5},
  frame           = single,
  framerule       = 0.4pt,
  rulecolor       = \color{codeFrame},
  breaklines      = true,
  showstringspaces= false,
  xleftmargin     = 1.5em,
  xrightmargin    = 0.5em,
  aboveskip       = 0.6em,
  belowskip       = 0.6em,
}

\lstdefinestyle{textstyle}{
  basicstyle      = \ttfamily\footnotesize,
  backgroundcolor = \color{black!4},
  frame           = single,
  framerule       = 0.4pt,
  rulecolor       = \color{codeFrame},
  breaklines      = true,
  xleftmargin     = 1.5em,
  xrightmargin    = 0.5em,
  aboveskip       = 0.6em,
  belowskip       = 0.6em,
}

% ── Callout boxes ──────────────────────────────────────────────────────────────
\tcbuselibrary{skins, breakable}

\newtcolorbox{tipbox}[1][]{
  enhanced, breakable,
  colback=tipteal!8, colframe=tipteal, coltitle=white,
  fonttitle=\bfseries\small, title={\faLightbulb\quad Tip#1},
  left=6pt, right=6pt, top=4pt, bottom=4pt,
  boxrule=0.8pt, arc=3pt,
}

\newtcolorbox{warnbox}[1][]{
  enhanced, breakable,
  colback=warnAmber!8, colframe=kmutnbGold, coltitle=white,
  fonttitle=\bfseries\small, title={\faExclamationTriangle\quad Note#1},
  left=6pt, right=6pt, top=4pt, bottom=4pt,
  boxrule=0.8pt, arc=3pt,
}

\newtcolorbox{checkpointbox}[1][]{
  enhanced,
  colback=kmutnbBlue!6, colframe=kmutnbBlue, coltitle=white,
  fonttitle=\bfseries\small, title={\faCheckCircle\quad Checkpoint#1},
  left=6pt, right=6pt, top=4pt, bottom=4pt,
  boxrule=0.8pt, arc=3pt,
}

\newtcolorbox{questionbox}[1][]{
  enhanced, breakable,
  colback=kmutnbGold!10, colframe=kmutnbGold, coltitle=white,
  fonttitle=\bfseries\small, title={\faQuestion\quad Question#1},
  left=6pt, right=6pt, top=4pt, bottom=4pt,
  boxrule=0.8pt, arc=3pt,
}

% ── Checkbox list ──────────────────────────────────────────────────────────────
\newlist{checklist}{itemize}{1}
\setlist[checklist]{label=\faSquare[regular], leftmargin=1.5em, noitemsep}

% ── Misc helpers ───────────────────────────────────────────────────────────────
\newcommand{\file}[1]{\texttt{\small #1}}
\newcommand{\api}[1]{\texttt{\small #1}}

% ── Title block macro ─────────────────────────────────────────────────────────
% Usage: \labtitle{03}{Aperiodic Servers \& Producer--Consumer}{2}{CLO 1 \& CLO 2}
\newcommand{\labtitle}[4]{%
  \begin{center}
    {\color{kmutnbBlue}\rule{\linewidth}{2pt}}\\[6pt]
    {\LARGE\bfseries\color{kmutnbBlue} Laboratory #1}\\[4pt]
    {\large\bfseries #2}\\[2pt]
    {\normalsize Real-Time Operating Systems \enspace\textbullet\enspace
     M.Eng.\ ECE \enspace\textbullet\enspace KMUTNB}\\[8pt]
    {\color{kmutnbGold}\rule{\linewidth}{1pt}}
  \end{center}
  \vspace{0.3cm}
  \begin{tabular}{@{}p{3.8cm} p{10cm}@{}}
    \textbf{Platform}       & NXP FRDM-MCXN236 (ARM Cortex-M33 @ 150\,MHz) \\[2pt]
    \textbf{RTOS}           & FreeRTOS v10.6 \\[2pt]
    \textbf{Toolchain}      & ARM GNU Toolchain (\texttt{arm-none-eabi-gcc}) + Make \\[2pt]
    \textbf{Estimated time} & #3 hours \\[2pt]
    \textbf{CLO mapping}    & #4 \\
  \end{tabular}
  \vspace{0.4cm}
}
  (from each labNN/ directory)
% ──────────────────────────────────────────────────────────────────────────────

% ── Packages ──────────────────────────────────────────────────────────────────
\usepackage[a4paper, top=2.5cm, bottom=2.5cm, left=2.5cm, right=2.5cm, headheight=28pt]{geometry}
\usepackage[utf8]{inputenc}
\usepackage[T1]{fontenc}
\usepackage{lmodern}
\usepackage{booktabs}
\usepackage{array}
\usepackage{tabularx}
\usepackage{longtable}
\usepackage{multirow}
\usepackage{colortbl}
\usepackage{xcolor}
\usepackage{titlesec}
\usepackage{enumitem}
\usepackage{hyperref}
\usepackage{fancyhdr}
\usepackage{parskip}
\usepackage{listings}
\usepackage[most]{tcolorbox}
\usepackage{fontawesome5}
\usepackage{microtype}
\usepackage{graphicx}
\usepackage{amsmath}

% ── Colors (KMUTNB brand) ──────────────────────────────────────────────────────
\definecolor{kmutnbBlue}{RGB}{0,56,116}
\definecolor{kmutnbGold}{RGB}{186,146,36}
\definecolor{lightGray}{RGB}{245,245,245}
\definecolor{midGray}{RGB}{200,200,200}
\definecolor{codeGray}{RGB}{248,248,248}
\definecolor{codeFrame}{RGB}{220,220,220}
\definecolor{tipteal}{RGB}{0,105,92}
\definecolor{warnAmber}{RGB}{121,85,0}
\definecolor{checkGreen}{RGB}{27,94,32}

% ── Section formatting ─────────────────────────────────────────────────────────
\titleformat{\section}
  {\large\bfseries\color{kmutnbBlue}}
  {\thesection.}{0.5em}{}
  [\color{kmutnbGold}\titlerule]

\titleformat{\subsection}
  {\normalsize\bfseries\color{kmutnbBlue}}
  {\thesubsection.}{0.5em}{}

\titleformat{\subsubsection}
  {\small\bfseries\color{kmutnbBlue}}
  {\thesubsubsection.}{0.5em}{}

\titlespacing*{\section}{0pt}{1.4ex plus .2ex}{0.6ex plus .1ex}
\titlespacing*{\subsection}{0pt}{1.0ex plus .2ex}{0.4ex plus .1ex}

% ── Listings (C and bash) ──────────────────────────────────────────────────────
\lstdefinestyle{cstyle}{
  language        = C,
  basicstyle      = \ttfamily\small,
  keywordstyle    = \bfseries\color{kmutnbBlue},
  commentstyle    = \itshape\color{gray},
  stringstyle     = \color{tipteal},
  numberstyle     = \tiny\color{midGray},
  numbers         = left,
  numbersep       = 6pt,
  stepnumber      = 1,
  backgroundcolor = \color{codeGray},
  frame           = single,
  framerule       = 0.4pt,
  rulecolor       = \color{codeFrame},
  breaklines      = true,
  showstringspaces= false,
  tabsize         = 4,
  xleftmargin     = 1.5em,
  xrightmargin    = 0.5em,
  aboveskip       = 0.8em,
  belowskip       = 0.8em,
}

\lstdefinestyle{bashstyle}{
  language        = bash,
  basicstyle      = \ttfamily\small,
  keywordstyle    = \color{kmutnbBlue},
  commentstyle    = \itshape\color{gray},
  backgroundcolor = \color{black!5},
  frame           = single,
  framerule       = 0.4pt,
  rulecolor       = \color{codeFrame},
  breaklines      = true,
  showstringspaces= false,
  xleftmargin     = 1.5em,
  xrightmargin    = 0.5em,
  aboveskip       = 0.6em,
  belowskip       = 0.6em,
}

\lstdefinestyle{textstyle}{
  basicstyle      = \ttfamily\footnotesize,
  backgroundcolor = \color{black!4},
  frame           = single,
  framerule       = 0.4pt,
  rulecolor       = \color{codeFrame},
  breaklines      = true,
  xleftmargin     = 1.5em,
  xrightmargin    = 0.5em,
  aboveskip       = 0.6em,
  belowskip       = 0.6em,
}

% ── Callout boxes ──────────────────────────────────────────────────────────────
\tcbuselibrary{skins, breakable}

\newtcolorbox{tipbox}[1][]{
  enhanced, breakable,
  colback=tipteal!8, colframe=tipteal, coltitle=white,
  fonttitle=\bfseries\small, title={\faLightbulb\quad Tip#1},
  left=6pt, right=6pt, top=4pt, bottom=4pt,
  boxrule=0.8pt, arc=3pt,
}

\newtcolorbox{warnbox}[1][]{
  enhanced, breakable,
  colback=warnAmber!8, colframe=kmutnbGold, coltitle=white,
  fonttitle=\bfseries\small, title={\faExclamationTriangle\quad Note#1},
  left=6pt, right=6pt, top=4pt, bottom=4pt,
  boxrule=0.8pt, arc=3pt,
}

\newtcolorbox{checkpointbox}[1][]{
  enhanced,
  colback=kmutnbBlue!6, colframe=kmutnbBlue, coltitle=white,
  fonttitle=\bfseries\small, title={\faCheckCircle\quad Checkpoint#1},
  left=6pt, right=6pt, top=4pt, bottom=4pt,
  boxrule=0.8pt, arc=3pt,
}

\newtcolorbox{questionbox}[1][]{
  enhanced, breakable,
  colback=kmutnbGold!10, colframe=kmutnbGold, coltitle=white,
  fonttitle=\bfseries\small, title={\faQuestion\quad Question#1},
  left=6pt, right=6pt, top=4pt, bottom=4pt,
  boxrule=0.8pt, arc=3pt,
}

% ── Checkbox list ──────────────────────────────────────────────────────────────
\newlist{checklist}{itemize}{1}
\setlist[checklist]{label=\faSquare[regular], leftmargin=1.5em, noitemsep}

% ── Misc helpers ───────────────────────────────────────────────────────────────
\newcommand{\file}[1]{\texttt{\small #1}}
\newcommand{\api}[1]{\texttt{\small #1}}

% ── Title block macro ─────────────────────────────────────────────────────────
% Usage: \labtitle{03}{Aperiodic Servers \& Producer--Consumer}{2}{CLO 1 \& CLO 2}
\newcommand{\labtitle}[4]{%
  \begin{center}
    {\color{kmutnbBlue}\rule{\linewidth}{2pt}}\\[6pt]
    {\LARGE\bfseries\color{kmutnbBlue} Laboratory #1}\\[4pt]
    {\large\bfseries #2}\\[2pt]
    {\normalsize Real-Time Operating Systems \enspace\textbullet\enspace
     M.Eng.\ ECE \enspace\textbullet\enspace KMUTNB}\\[8pt]
    {\color{kmutnbGold}\rule{\linewidth}{1pt}}
  \end{center}
  \vspace{0.3cm}
  \begin{tabular}{@{}p{3.8cm} p{10cm}@{}}
    \textbf{Platform}       & NXP FRDM-MCXN236 (ARM Cortex-M33 @ 150\,MHz) \\[2pt]
    \textbf{RTOS}           & FreeRTOS v10.6 \\[2pt]
    \textbf{Toolchain}      & ARM GNU Toolchain (\texttt{arm-none-eabi-gcc}) + Make \\[2pt]
    \textbf{Estimated time} & #3 hours \\[2pt]
    \textbf{CLO mapping}    & #4 \\
  \end{tabular}
  \vspace{0.4cm}
}


% ── Header / Footer ────────────────────────────────────────────────────────────
\pagestyle{fancy}
\fancyhf{}
\fancyhead[L]{\small\color{kmutnbBlue}\textbf{KMUTNB -- Faculty of Engineering}}
\fancyhead[R]{\small\color{kmutnbBlue}Dept.\ of Electrical and Computer Engineering}
\fancyfoot[L]{\small\color{kmutnbBlue}Real-Time Operating Systems -- M.Eng.\ ECE}
\fancyfoot[C]{\small\thepage}
\fancyfoot[R]{\small\color{kmutnbBlue}Lab 11}
\renewcommand{\headrulewidth}{0.4pt}
\renewcommand{\headrule}{\hbox to\headwidth{%
  \color{kmutnbGold}\leaders\hrule height\headrulewidth\hfill}}
\renewcommand{\footrulewidth}{0.4pt}
\renewcommand{\footrule}{\hbox to\headwidth{%
  \color{midGray}\leaders\hrule height\footrulewidth\hfill}}

\hypersetup{
  colorlinks = true,
  linkcolor  = kmutnbBlue,
  urlcolor   = kmutnbBlue,
  citecolor  = kmutnbBlue,
  pdftitle   = {Lab 11 -- Zephyr RTOS: Port, Device Tree and Kconfig},
  pdfauthor  = {KMUTNB RTOS Course},
}

% ══════════════════════════════════════════════════════════════════════════════
\begin{document}

\labtitle{11}{Zephyr RTOS: Port, Device Tree \& Kconfig}{4--5}{CLO 5 \& CLO 6}

% ── Objectives ────────────────────────────────────────────────────────────────
\section{Objectives}

By completing this lab you will be able to:

\begin{enumerate}[noitemsep, leftmargin=*]
  \item \textbf{Set up} a Zephyr RTOS workspace for the FRDM-MCXN236 using \texttt{west}.
  \item \textbf{Port} the Lab 03 producer-consumer pipeline from FreeRTOS to Zephyr
        threading APIs.
  \item \textbf{Write} a Device Tree overlay to add a simulated BME280 sensor on I2C.
  \item \textbf{Use} the Zephyr Sensor API (\api{sensor\_sample\_fetch} /
        \api{sensor\_channel\_get}) with a DT-bound driver.
  \item \textbf{Compare} binary size, API surface, and configuration overhead between
        FreeRTOS and Zephyr.
\end{enumerate}

% ── Prerequisites ─────────────────────────────────────────────────────────────
\section{Prerequisites}

\begin{checklist}
  \item Lab 03 producer-consumer source (for reference during porting).
  \item Python 3.8+ and \texttt{pip} available.
  \item At least 4\,GB free disk space (Zephyr workspace + HALs).
  \item Git installed.
  \item USB connection to FRDM-MCXN236 (J-Link or MCU-Link).
\end{checklist}

\begin{warnbox}
  \textbf{Time note:} \texttt{west update} downloads $\sim$1.5\,GB of source. Do this
  on a fast network connection. The download is a one-time setup; subsequent builds are
  fast.
\end{warnbox}

% ── Background ────────────────────────────────────────────────────────────────
\section{Background: Zephyr vs.\ FreeRTOS}

\renewcommand{\arraystretch}{1.3}
\begin{tabular}{@{} p{3.2cm} p{5.2cm} p{4.8cm} @{}}
  \toprule
  \textbf{Aspect} & \textbf{FreeRTOS} & \textbf{Zephyr} \\
  \midrule
  Configuration  & \file{FreeRTOSConfig.h} (\texttt{\#define}) & \file{prj.conf} (Kconfig) \\
  Hardware desc. & Code (\texttt{GPIO\_PinInit(...)}) & Device Tree (\file{.dts} / \file{.overlay}) \\
  Build system   & Make / CMake (manual) & \texttt{west} (unified) \\
  Driver model   & SDK-specific (\file{fsl\_gpio.h}) & Unified API (\file{gpio.h}) \\
  Networking/BT  & External libraries & Built-in subsystems \\
  Certification  & Qualified kits available & Partial (Safety Cert.\ project) \\
  \bottomrule
\end{tabular}
\renewcommand{\arraystretch}{1}

% ── Project Structure ─────────────────────────────────────────────────────────
\section{Project Structure}

\begin{lstlisting}[style=textstyle]
lab11_zephyr/
  CMakeLists.txt
  prj.conf                    <- Kconfig (replaces FreeRTOSConfig.h)
  app.overlay                 <- Device Tree customisation
  src/
    main.c                    <- Zephyr entry, LOG_INF banner
    sensor_thread.c/.h        <- producer (replaces vSensor1Task)
    logger_thread.c/.h        <- consumer (replaces vLoggerTask)
  boards/
    frdm_mcxn236.conf         <- board-specific Kconfig overrides (optional)
\end{lstlisting}

% ── Part A ─────────────────────────────────────────────────────────────────────
\section{Part A --- Workspace Setup}

\subsection{Step 1 --- Install west}

\begin{lstlisting}[style=bashstyle]
pip install west
\end{lstlisting}

\subsection{Step 2 --- Initialise the Zephyr Workspace}

\begin{lstlisting}[style=bashstyle]
mkdir ~/zephyrproject && cd ~/zephyrproject
west init -m https://github.com/zephyrproject-rtos/zephyr --mr v3.6.0
west update          # ~20 minutes on first run
west zephyr-export   # registers CMake package
pip install -r zephyr/scripts/requirements.txt
\end{lstlisting}

\subsection{Step 3 --- Verify with hello\_world}

\begin{lstlisting}[style=bashstyle]
cd ~/zephyrproject
west build -b frdm_mcxn236 zephyr/samples/hello_world
west flash
\end{lstlisting}

Expected terminal output:

\begin{lstlisting}[style=textstyle]
*** Booting Zephyr OS build v3.6.0 ***
Hello World! frdm_mcxn236
\end{lstlisting}

\begin{checkpointbox}[~A]
  Screenshot of \texttt{hello\_world} running on the board.
\end{checkpointbox}

\begin{questionbox}[~A1]
  \texttt{west update} fetches many modules (HALs, mbedTLS, LVGL, etc.). Which module
  provides the NXP MCX HAL for the FRDM-MCXN236? (Hint: look in
  \texttt{zephyrproject/modules/hal/}.)
\end{questionbox}

% ── Part B ─────────────────────────────────────────────────────────────────────
\section{Part B --- Port the Producer-Consumer Pipeline}

\subsection{Step 1 --- FreeRTOS $\to$ Zephyr API Mapping}

Complete this mapping table before writing any code:

\renewcommand{\arraystretch}{1.3}
\begin{tabular}{@{} p{6.5cm} p{6cm} @{}}
  \toprule
  \textbf{FreeRTOS (Lab 03)} & \textbf{Zephyr equivalent} \\
  \midrule
  \api{xTaskCreate(fn, name, stack\_words, ...)} & \\
  \api{K\_THREAD\_DEFINE(tid, stack\_bytes, fn, ...)} & \\
  \api{xQueueCreate(depth, item\_size)} & \\
  \api{xQueueSend(q, \&item, timeout)} & \\
  \api{xQueueReceive(q, \&item, portMAX\_DELAY)} & \\
  \api{xSemaphoreCreateBinary()} & \\
  \api{xSemaphoreGive(sem)} & \\
  \api{xSemaphoreTake(sem, portMAX\_DELAY)} & \\
  \api{vTaskDelay(pdMS\_TO\_TICKS(n))} & \\
  \api{vTaskDelayUntil(\&last, period)} & \\
  \bottomrule
\end{tabular}
\renewcommand{\arraystretch}{1}

\subsection{Step 2 --- \texttt{prj.conf}}

\begin{lstlisting}[style=textstyle]
# prj.conf
CONFIG_SERIAL=y
CONFIG_LOG=y
CONFIG_LOG_DEFAULT_LEVEL=3

# Threading
CONFIG_MAIN_STACK_SIZE=2048
CONFIG_HEAP_MEM_POOL_SIZE=4096

# No float printf -- save flash
CONFIG_CBPRINTF_FP_SUPPORT=n
\end{lstlisting}

\subsection{Step 3 --- \texttt{CMakeLists.txt}}

\begin{lstlisting}[style=textstyle]
cmake_minimum_required(VERSION 3.20)
find_package(Zephyr REQUIRED HINTS $ENV{ZEPHYR_BASE})
project(lab11_zephyr)

target_sources(app PRIVATE
    src/main.c
    src/sensor_thread.c
    src/logger_thread.c
)
\end{lstlisting}

\subsection{Step 4 --- Sensor Producer Thread}

In \file{src/sensor\_thread.c}:

\begin{lstlisting}[style=cstyle]
#include <zephyr/kernel.h>
#include "sensor_thread.h"

/* Message queue: 10 items of SensorReading_t */
K_MSGQ_DEFINE(sensor_msgq, sizeof(SensorReading_t), 10, 4);

static void sensor_thread_fn(void *a, void *b, void *c)
{
    SensorReading_t reading = { .sensor_id = 1, .value = 0 };

    while (1) {
        k_sleep(K_MSEC(100));
        reading.value = (reading.value + 7) % 1000;
        reading.ts_ms = k_uptime_get_32();

        if (k_msgq_put(&sensor_msgq, &reading, K_NO_WAIT) != 0)
            LOG_WRN("queue full -- dropped");
    }
}

K_THREAD_DEFINE(sensor_tid, 1024, sensor_thread_fn,
                NULL, NULL, NULL, 5, 0, 0);
\end{lstlisting}

\subsection{Step 5 --- Logger Consumer Thread}

In \file{src/logger\_thread.c}:

\begin{lstlisting}[style=cstyle]
#include <zephyr/kernel.h>
#include <zephyr/logging/log.h>
#include "sensor_thread.h"   /* for sensor_msgq extern */

LOG_MODULE_REGISTER(logger, LOG_LEVEL_INF);

static void logger_thread_fn(void *a, void *b, void *c)
{
    SensorReading_t r;
    uint32_t count = 0;
    while (1) {
        k_msgq_get(&sensor_msgq, &r, K_FOREVER);
        count++;
        LOG_INF("#%u  sensor=%d  val=%d  t=%u ms",
                count, r.sensor_id, r.value, r.ts_ms);
    }
}

K_THREAD_DEFINE(logger_tid, 1024, logger_thread_fn,
                NULL, NULL, NULL, 3, 0, 0);
\end{lstlisting}

\subsection{Step 6 --- Build and Flash}

\begin{lstlisting}[style=bashstyle]
cd ~/zephyrproject/lab11_zephyr
west build -b frdm_mcxn236 .
west flash
\end{lstlisting}

Expected UART output:

\begin{lstlisting}[style=textstyle]
[00:00:00.012,000] <inf> logger: #1  sensor=1  val=7  t=100 ms
[00:00:00.112,000] <inf> logger: #2  sensor=1  val=14  t=200 ms
...
\end{lstlisting}

\begin{checkpointbox}[~B]
  Terminal showing Zephyr logging output from the ported pipeline.
\end{checkpointbox}

\begin{questionbox}[~B1]
  In FreeRTOS, \api{xQueueSend} with \api{portMAX\_DELAY} blocks the calling task
  forever. What is the Zephyr equivalent timeout value for ``wait forever''?
\end{questionbox}

\begin{questionbox}[~B2]
  \api{K\_THREAD\_DEFINE} takes stack size in \textbf{bytes} while FreeRTOS
  \api{xTaskCreate} takes it in \textbf{words} (4 bytes each). Convert the Lab 03
  sensor task's 128-word stack to the correct Zephyr byte value.
\end{questionbox}

% ── Part C ─────────────────────────────────────────────────────────────────────
\section{Part C --- Device Tree Sensor (BME280)}

\subsection{Step 1 --- Write the Overlay}

Create \file{app.overlay} in the project root:

\begin{lstlisting}[style=textstyle]
/* app.overlay -- add BME280 sensor on LPI2C1 */
&lpi2c1 {
    status = "okay";
    clock-frequency = <I2C_BITRATE_STANDARD>;
    pinctrl-0 = <&pinmux_lpi2c1>;
    pinctrl-names = "default";

    bme280: bme280@76 {
        compatible = "bosch,bme280";
        reg = <0x76>;
        status = "okay";
    };
};
\end{lstlisting}

Add to \file{prj.conf}:

\begin{lstlisting}[style=textstyle]
CONFIG_I2C=y
CONFIG_SENSOR=y
CONFIG_BME280=y
\end{lstlisting}

\subsection{Step 2 --- Use the Sensor API}

Add a BME280 reading thread in \file{src/sensor\_thread.c}:

\begin{lstlisting}[style=cstyle]
#include <zephyr/drivers/sensor.h>

static const struct device *bme280 =
    DEVICE_DT_GET(DT_NODELABEL(bme280));

static void bme_thread_fn(void *a, void *b, void *c)
{
    if (!device_is_ready(bme280)) {
        LOG_ERR("BME280 not ready");
        return;
    }

    struct sensor_value temp, pressure, humidity;
    while (1) {
        k_sleep(K_MSEC(1000));
        sensor_sample_fetch(bme280);
        sensor_channel_get(bme280, SENSOR_CHAN_AMBIENT_TEMP, &temp);
        sensor_channel_get(bme280, SENSOR_CHAN_PRESS, &pressure);
        sensor_channel_get(bme280, SENSOR_CHAN_HUMIDITY, &humidity);
        LOG_INF("T=%.2f C  P=%.2f hPa  H=%.2f %%",
                sensor_value_to_double(&temp),
                sensor_value_to_double(&pressure) / 1000.0,
                sensor_value_to_double(&humidity));
    }
}

K_THREAD_DEFINE(bme_tid, 1024, bme_thread_fn, NULL, NULL, NULL, 4, 0, 0);
\end{lstlisting}

\begin{warnbox}
  If you do not have a physical BME280, use the Zephyr emulator by adding to
  \file{prj.conf}:
  \begin{lstlisting}[style=textstyle, aboveskip=2pt, belowskip=0pt]
CONFIG_BME280_TRIGGER_NONE=y
CONFIG_EMUL=y
CONFIG_BME280_EMUL=y
  \end{lstlisting}
  The emulator returns synthetic values, letting you exercise the full API without
  hardware.
\end{warnbox}

\begin{checkpointbox}[~C]
  Terminal showing BME280 (or emulated) temperature/pressure/humidity readings every
  second.
\end{checkpointbox}

\begin{questionbox}[~C1]
  Where in the Zephyr source tree is the BME280 driver located? What is the value of
  its \texttt{compatible} string in the DTS binding file?
\end{questionbox}

\begin{questionbox}[~C2]
  You want to add a second BME280 at I2C address 0x77 on the same bus. What change do
  you make --- in \file{app.overlay} only, or also in C code? Why?
\end{questionbox}

% ── Part D ─────────────────────────────────────────────────────────────────────
\section{Part D --- Binary Size Comparison}

\subsection{Step 1 --- Measure Zephyr Binary Size}

\begin{lstlisting}[style=bashstyle]
west build -b frdm_mcxn236 .
arm-none-eabi-size build/zephyr/zephyr.elf
\end{lstlisting}

\subsection{Step 2 --- Measure Equivalent FreeRTOS Binary Size}

\begin{lstlisting}[style=bashstyle]
cd ../lab03_aperiodic_server
make
arm-none-eabi-size lab03.elf
\end{lstlisting}

\subsection{Step 3 --- Complete the Comparison Table}

\renewcommand{\arraystretch}{1.3}
\begin{tabular}{@{} p{4cm} p{3.5cm} p{3.5cm} @{}}
  \toprule
  \textbf{Metric} & \textbf{FreeRTOS (Lab 03)} & \textbf{Zephyr (Lab 11)} \\
  \midrule
  \texttt{.text} (code)          & & \\
  \texttt{.rodata} (const)       & & \\
  \texttt{.data} + \texttt{.bss} & & \\
  Total flash                    & & \\
  Build time (cold)              & & \\
  Lines of config                & & \\
  \bottomrule
\end{tabular}
\renewcommand{\arraystretch}{1}

\begin{questionbox}[~D1]
  Zephyr is typically larger than FreeRTOS for the same feature set. Where does most
  of the extra flash usage come from? (Hint: run \texttt{west build -t rom\_report}.)
\end{questionbox}

\begin{questionbox}[~D2]
  You need to add BLE advertising to this application. Describe how you would do this
  in Zephyr (which \file{prj.conf} options, which API) vs.\ how you would approach it
  with FreeRTOS (which external library, which integration steps). Which is faster to
  integrate and why?
\end{questionbox}

% ── Part E (Optional) ─────────────────────────────────────────────────────────
\section{Part E --- Kconfig Exploration (Optional Challenge)}

\begin{lstlisting}[style=bashstyle]
west build -b frdm_mcxn236 -t menuconfig .
\end{lstlisting}

Navigate to:
\begin{itemize}[noitemsep, leftmargin=*]
  \item \texttt{Kernel $\to$ Threads} --- find \texttt{CONFIG\_MAIN\_STACK\_SIZE}
  \item \texttt{Subsystems $\to$ Logging} --- observe log level options and backend
        choices
  \item \texttt{Drivers $\to$ Sensor} --- note available sensor drivers
\end{itemize}

\begin{questionbox}[~Challenge]
  Disable the shell (\texttt{CONFIG\_SHELL=n}) and the logging backend to serial
  (\texttt{CONFIG\_LOG\_BACKEND\_UART=n}). Rebuild and measure the flash savings. How
  much code do logging and the shell add to a Zephyr application?
\end{questionbox}

% ── Deliverables ───────────────────────────────────────────────────────────────
\section{Deliverables}

\renewcommand{\arraystretch}{1.3}
\begin{tabular}{@{} c p{5cm} p{6.5cm} @{}}
  \toprule
  \textbf{\#} & \textbf{Item} & \textbf{Details} \\
  \midrule
  1 & Part A screenshot & \texttt{hello\_world} running on the board \\
  2 & Part B API mapping table & Completed FreeRTOS $\to$ Zephyr column \\
  3 & Part B terminal & Zephyr producer-consumer output \\
  4 & Part C terminal & BME280 (real or emulated) sensor readings \\
  5 & Part D table & Binary size comparison \\
  6 & File listings & \file{prj.conf} and \file{app.overlay} \\
  7 & Written answers & Questions A1, B1--B2, C1--C2, D1--D2 \\
  8 & Part E (optional) & Flash savings measurement \\
  \bottomrule
\end{tabular}
\renewcommand{\arraystretch}{1}

% ── Key Concepts ───────────────────────────────────────────────────────────────
\section{Key Concepts Demonstrated}

\renewcommand{\arraystretch}{1.3}
\begin{tabular}{@{} p{6.5cm} p{6cm} @{}}
  \toprule
  \textbf{Concept} & \textbf{Where you saw it} \\
  \midrule
  west workspace init and build & Part A \\
  Kconfig compile-time configuration & Part B \file{prj.conf} \\
  \api{K\_THREAD\_DEFINE} and \api{K\_MSGQ\_DEFINE} & Part B porting \\
  Device Tree overlay for I2C sensor & Part C \\
  Unified Sensor API & Part C \\
  FreeRTOS vs.\ Zephyr binary size & Part D \\
  \bottomrule
\end{tabular}
\renewcommand{\arraystretch}{1}

% ── Reference ─────────────────────────────────────────────────────────────────
\section{Reference}

\renewcommand{\arraystretch}{1.3}
\begin{tabular}{@{} p{5cm} p{7.5cm} @{}}
  \toprule
  \textbf{Resource} & \textbf{Relevant sections} \\
  \midrule
  Zephyr documentation & Getting Started, Threads, Message Queue, Sensor API \\
  Zephyr Device Tree & \file{boards/arm/frdm\_mcxn236/},
                       \file{dts/bindings/sensor/bosch,bme280.yaml} \\
  Nordic DevAcademy & NCS Fundamentals --- Device Tree module \\
  FreeRTOS docs & \api{xQueueCreate} (for comparison) \\
  \bottomrule
\end{tabular}
\renewcommand{\arraystretch}{1}

\end{document}
